\section{Fine-Tuning Setup}

The repository fine-tuning path is intentionally local-first. Training begins with expert highway trajectories collected under the same task family used for evaluation, then converts those trajectories into a strict instruction-output format suitable for supervised adaptation. The pipeline subsequently rebalances the action distribution, validates row schema and action formatting, produces merged model exports after training, and packages those exports for local deployment through GGUF conversion and Ollama model creation. This sequence matters for the paper because the evaluated fine-tuned models are not abstract checkpoints: they are deployment-oriented local artifacts produced by the same toolchain that a practitioner would use to run the benchmark end to end.

The training implementation is compatible with efficient adaptation in the QLoRA line \cite{dettmers2023qlora}, but the paper does not claim novelty in the optimizer or adapter mechanism itself. The methodological question is instead whether locally adapted models remain meaningfully comparable once they are placed back under the same runtime-constrained driving benchmark. For that reason, the fine-tuning section is primarily about comparison control rather than optimization design.

To make that comparison explicit, the study organizes models through a registry that records family, tier, \texttt{variant\_kind}, and a \texttt{base\_model\_id} link from each fine-tuned artifact back to its intended base model. This registry provides nominal lineage and lets the post-hoc study pipeline recover base-versus-fine-tuned pair structure after benchmark execution. Lineage alone, however, is not enough for scientific comparison. We make pairwise fine-tuning claims only when both sides of a family remain ranking-eligible under the same benchmark regime. This rule is a methodological control, not merely a reporting preference: if either the base model or the adapted model falls out of the valid ranking set, then a base-versus-tuned delta can be driven by timeout collapse or invalid execution rather than adaptation quality.

Under the current publication bundle, the only exact eligible pair is the lightweight Phi family. The paper should therefore be read as an initial fine-tuning study with narrow but real causal evidence, not as a complete family-by-family adaptation benchmark. Lightweight scale still supports one direct before-and-after comparison under shared runtime conditions. Midclass and highclass, by contrast, still contain training artifacts and registry-defined lineages, but they do not preserve the benchmark conditions needed for equally strong pairwise interpretation. That boundary is central to the paper's methodology: we compare only where the benchmark still behaves like a benchmark, and we treat the loss of pairwise comparability as part of the empirical result rather than as a missing-data nuisance.
